\section{Conclusion}\label{sec:conclusion}

In this work, we implemented a bitwise version of the Metropolis algorithm which we validate on a number of models including the Hopfield model. This version uses the machine-specific binary arithmetic to efficiently simulate the evolution of $n_\text{bits}$ systems in parallel. On top of this structural optimization, the same sequence of random numbers is used to make all $n_\text{bits}$ systems evolve. This turns out to be central as random number generation is often the simulation bottleneck. Using this bitwise formalism considerably reduces the number of operations required to make the system evolve, as all configurations are stored using \Bits objects, whose evolutions are given by bitwise operations that are straightforward for the machine. For now, we use $n_\text{bits}=64$, but this number can potentially be increased up $n_\text{bits}=512$, providing a change in the type of objects used and low-level implementation of the code. The limitations of the size of $n_\text{bits}$ are directly related to the constraints of the CPU architecture.\\
For now, the code only provides a bitwise implementation of the Metropolis algorithm, but other Monte-Carlo algorithms such as the Glauber algorithm \cite{Glauber1963} or cluster flip algorithms \cite{SwendsenWang1987,Wolff1989} can also be setup in this formalism. So far, the code enables the application of this bitwise Metropolis algorithm on several networks models: the Ising model, the Edwards-Anderson spin glass model and the Hopfield model. It can however be generalized to every model that uses discrete values, as it is always possible to shift those values into positive integers that are suitably stored in \UnsignedInt objects.\\
The interest of this bitwise formalism is to efficiently sample the dynamics of a given system. Indeed, the Metropolis algorithm becomes computationally demanding when one attempts to sample a large number of configurations. This situation occurs, for example, when looking for rare samples in the dynamics, or sampling a large number of initial conditions.\\
This implementation is therefore well suited to studying the dynamics of the Hopfield model. While the equilibrium properties of the model -- introduced in its now-canonical form by Hopfield in 1982 \cite{Hopfield1982} -- are well understood, its dynamics has been solved analytically only in limiting cases, such as strongly diluted networks \cite{Derrida1987} or the thermodynamic limit \cite{Coolen2001}, and numerical simulations remain essential to study finite networks, rare events or the statistics over many initial conditions. \\
The code has first been tested on the Ising and Edwards-Anderson models, for which it reproduced the correct expected results, see \Cref{fig:phase_diagram_ising} and \Cref{fig:binder}. In addition, it has proven to be efficient, yielding computational gains -- that is the ratio of the number of independent results obtained for a given computation time between the optimized bitwise implementation and the reference one -- comprised between 20 and 45, depending mostly on the temperature of the system. However, when applying the same bitwise implementation on the Hopfield network, the computational gain vanished. This is most likely due to the density of the network: one has to go through every single neuron in the network to update a single neuron. If the bitwise implementation creates some temporary \Bits or \UnsignedInt objects during the spin interaction calculations, the loss in performance can become considerable for dense networks. Despite this loss of computational advantage, the bitwise implementation still yielded the expected results, see \Cref{fig:reconstruction_error}. In addition, this implementation allows us to easily visualize the evolution of the $n_\text{bits}$ independent system in parallel, see \Cref{fig:evolution_ising}, \Cref{fig:all_mu_single_replica} and \Cref{fig:mu_star_all_replica}.\\
This method is thus a promising tool to study the dynamics of the Hopfield network, such as its dynamical exponents, the correlations between configurations, or the dependence of the dynamics on the structure of the network.
